% Pista alg-26j2-q2 · Àlgebra
% Procedència: PAU Catalunya 2026 · Convocatòria de juny · Sèrie 5 · Exercici 2
\documentclass[11pt,a4paper]{article}
\usepackage{pista}
\pistainfo{Catalunya 2026}{Convocatòria de juny}{Sèrie 5}

\begin{document}

\section*{\textcolor{blauPista}{Exercici 2}}

\noindent Considereu les matrius
\[
M = \begin{pmatrix} 3 & 3 & 0 \\ k - 1 & -1 & -1 \end{pmatrix} \quad \text{i} \quad
N = \begin{pmatrix} 3 & 1 \\ k - 1 & 0 \\ 0 & 2 \end{pmatrix},
\]
on $k$ és un nombre real qualsevol.

\subsection*{2.1. \normalfont Calculeu els valors de $k$ per als quals $MN$ és invertible. \hspace*{\fill}\textnormal{[0,75~punts]}}

\pista{El producte $MN$ (una matriu $2 \times 3$ per una $3 \times 2$) és una matriu $2 \times 2$.
\begin{itemize}
\item Calcula'l explícitament: alguns elements depenen de $k$.
\item Calcula'n el determinant (t'ha de sortir una expressió de segon grau en $k$) i troba per a quins valors s'anul·la.
\item Per a aquests valors, $MN$ no és invertible; per a tots els altres, sí.
\end{itemize}}

\subsection*{2.2. \normalfont Calculeu els valors de $k$ per als quals $NM$ és invertible. \hspace*{\fill}\textnormal{[0,75~punts]}}

\pista{Ara $NM$ és una matriu $3 \times 3$: calcula el producte.

\medskip
Abans de desenvolupar el determinant, fixa't en les columnes:
\begin{itemize}
\item Alguna és combinació lineal de les altres dues? (Els coeficients poden dependre de $k$.)
\item Si és així, què en dedueixes sobre el determinant i, per tant, sobre la invertibilitat d'$NM$?
\end{itemize}}

\subsection*{2.3. \normalfont Estudieu per a quins valors reals de $m$, $n$ i $k$ el sistema $M \begin{pmatrix} x \\ y \\ z \end{pmatrix} = \begin{pmatrix} m \\ n \end{pmatrix}$ és compatible indeterminat. \hspace*{\fill}\textnormal{[1~punt]}}

\pista{Escriu les dues equacions del sistema: la matriu del sistema és $M$, i l'ampliada, $M^{'}$.
\begin{itemize}
\item Busca un menor d'ordre 2 d'$M$ que no s'anul·li per a cap valor de $k$.
\item Troba $\mathrm{rang}(M)$ i $\mathrm{rang}(M^{'})$.
\item Decideix per a quins valors d'$m$, $n$ i $k$ el sistema és compatible indeterminat.
\end{itemize}}

\end{document}
